% Appendix 4: Polar Form of Two-Dimensional Vectors
% QLD Specialist Mathematics — Unit 1
\begingroup
\small

\noindent\fbox{%
\begin{minipage}{0.97\textwidth}
\textit{\textbf{Enrichment only.}
Polar form of vectors is part of
\textbf{QLD Specialist Mathematics (Unit~1)}.
It is not prescribed in the current NESA Extension~1 or Extension~2
syllabus.}
\end{minipage}}

\vspace{0.8em}

\noindent\textbf{Key concepts.}

A vector in two dimensions has both magnitude and direction.
We usually express a vector in Cartesian form:
\[
  \mathbf{v} = x\,\mathbf{i} + y\,\mathbf{j}.
\]
An alternative is \textbf{polar form}, which represents a non-zero
vector by
\[
  \boxed{(r,\,\theta)}
\]
where:
\begin{itemize}[leftmargin=*,itemsep=2pt]
  \item $r > 0$ is the magnitude (length) of the vector.
  \item $\theta$ is its direction, measured anticlockwise from the
        positive $x$-axis.
\end{itemize}

The two forms are related by
\[
  \boxed{x = r\cos\theta, \qquad y = r\sin\theta}.
\]
Conversely,
\[
  r = \sqrt{x^2 + y^2},
  \qquad
  \tan\theta = \frac{y}{x} \quad (x \ne 0).
\]

When calculating $\theta$, always check the quadrant.
The direction is determined modulo $2\pi$.
The zero vector has no uniquely defined direction.

\vspace{0.4em}
\noindent\textit{NSW connection.}
Extension~1 already teaches vector components, magnitude and direction,
while Extension~2 uses trigonometric vector equations of circles.
The enrichment here is the systematic use of polar-form notation
and conversion.

\vspace{0.8em}
\noindent\textbf{Worked example.}
Express
\[
  \mathbf{v} = -\sqrt{3}\,\mathbf{i} + \mathbf{j}
\]
in polar form.

\vspace{0.3em}
\noindent\textit{Step~1: Find the magnitude.}
\[
  r = \sqrt{(-\sqrt{3})^{2} + 1^{2}} = 2.
\]

\noindent\textit{Step~2: Find the direction.}
Since $x < 0$ and $y > 0$, the vector lies in Quadrant~II.
\[
  \tan\theta = -\frac{1}{\sqrt{3}}.
\]
Taking $0 \le \theta < 2\pi$,
\[
  \theta = \frac{5\pi}{6}.
\]
Therefore,
\[
  \boxed{(r,\,\theta) = \left(2,\;\frac{5\pi}{6}\right)}.
\]

\noindent\textit{Verification.}
\[
  2\cos\frac{5\pi}{6} = -\sqrt{3},
  \qquad
  2\sin\frac{5\pi}{6} = 1.  \quad\checkmark
\]

\vspace{0.8em}
\noindent\textbf{Key takeaways.}
\begin{itemize}[leftmargin=*,itemsep=2pt]
  \item Cartesian form is convenient for adding vectors and performing
        algebra; polar form makes magnitude and direction immediately
        visible.
  \item Always check the quadrant when converting to polar form.
  \item Polar form connects naturally with the modulus and argument
        of complex numbers.
\end{itemize}

\vspace{0.4em}
\noindent\textit{Beyond school --- University connection.}
Polar coordinates are widely used in physics, robotics and engineering.
They also connect vector geometry with complex numbers through
\[
  z = r(\cos\theta + i\sin\theta).
\]

\endgroup
